\documentclass[11pt]{article}
\usepackage[a4paper,margin=28mm]{geometry}
\usepackage{amsmath,amssymb,amsthm}
\usepackage{booktabs}
\usepackage[shortlabels]{enumitem}
\usepackage[colorlinks=true,citecolor=black,linkcolor=black,urlcolor=blue]{hyperref}
\setlength{\emergencystretch}{2em}
\renewcommand{\thesection}{S\arabic{section}}
\renewcommand{\thetable}{S\arabic{table}}
\newtheorem{lemma}{Lemma}
\renewcommand{\thelemma}{S\arabic{lemma}}

\title{Supplement to ``The maximum size of a blackout that keeps every lattice rectangle
identifiable''}
\author{Arjun Pemmasani}
\date{}

\begin{document}
\maketitle

This supplement describes the SAT encodings behind the upper bounds of the theorem on small
heights, proves that they are faithful, and lists the instances, the certificates, and the
formal proofs. We refer to results of the paper by their names. Paths are relative to the
root of the repository \url{https://github.com/apemm/rectangle-blackouts}.

\section{The SAT encodings}

Let $C(R)$ be the corner set of a rectangle $R$. A blackout $S$ is valid if and only if it
contains no set $C(R_1)\triangle C(R_2)$ with $R_1\ne R_2$ (the lemma on presentations in
Section 2 of the paper).

\emph{Encoding A} has a variable for each grid point and a clause for each
inclusion-minimal set of the form $C(R_1)\triangle C(R_2)$ with $R_1\ne R_2$. The clause
says that this set is not fully black.

\emph{Encoding B} has such clauses only for the pairs of rectangles that share a corner.
It also has a variable $y_R$ for each rectangle $R$, which is forced to be true if $R$ is
fully black, and it allows at most 1 true $y_R$.

A \emph{structured} instance on 6 rows fixes one row graph from the equality structure
theorem (a $K_{3,3}$, for $m+9$ points) or from the lemma on 6 rows and $m+8$ points
($K_{3,3}$ minus an edge, $K_{2,4}$, or $K_{3,3}$ with 1 empty column). Each column chooses
a type, which is a single row, an edge of the graph, or the empty set in the last case. Each
edge is used once, and the empty type is used once in the last case. The instance keeps
those clauses of encoding B that involve a tilted rectangle and whose set can be fully
black for these types. The size of the blackout is fixed by a sequential
counter~\cite{sinz} or a totalizer~\cite{bailleux}.

\begin{lemma}\label{lem:enc}
Let $S$ be a blackout of the $h\times m$ grid.
\begin{enumerate}[(a)]
\item $S$ is valid if and only if it contains no inclusion-minimal set of the form
$C(R_1)\triangle C(R_2)$ with $R_1\ne R_2$.
\item $S$ is valid if and only if at most 1 rectangle is fully black and $S$ contains no
set $C(R_1)\triangle C(R_2)$ for distinct rectangles $R_1,R_2$ that share a corner.
\item Let $h=6$ and $m\ge3$, and let $S$ be valid with $m+9$ or $m+8$ points. Then for one
of the row graphs above, the choice of the black rows of column $i$ as the type of column
$i$ extends to a satisfying assignment of the structured instance.
\end{enumerate}
\end{lemma}

\begin{proof}
(a) Every set $C(R_1)\triangle C(R_2)$ contains an inclusion-minimal one. (b) For rectangles
without a common corner, the set $C(R_1)\triangle C(R_2)$ is $C(R_1)\cup C(R_2)$, which is
black exactly when both rectangles are fully black. (c) The equality structure theorem or
the lemma on $m+8$ points gives the row graph, and the black rows of every column form an
allowed type. Set $y_R$ to true exactly for the fully black rectangles. By (b) this
satisfies encoding B, and the structured instance has a subset of its clauses.
\end{proof}

Hence the point variables of encodings A and B describe exactly the valid blackouts of the
chosen size. On 6 rows, $f(6,m)\le m+8$ once the 10 structured instances with $m+9$ points
are unsatisfiable, and $f(6,m)\le m+7$ once the further $115$ instances with $m+8$ points
are.

\section{Instances and certificates}

CaDiCaL~\cite{cadical} produced an LRAT proof for every unsatisfiable instance in
Table~\ref{tab:instances}. Each of the $739$ proofs was accepted by 2 checkers, the program
\texttt{lrat-check} of drat-trim~\cite{drattrim,cruzfilipe} and the LRAT checker of the
Lean core library. The proofs are large and are not stored. Their hashes, the hashes of
the CNF files, and the results of both checkers are in
\nolinkurl{jis/data/unsat_certificates.json} and
\nolinkurl{jis/data/six_row_excess8_certificates.json}. The pipeline in
\nolinkurl{jis/unsat} recreates the proofs. Two separately written generators produce the
same clauses for the instances of encoding A and for the $70$ structured instances with
$m+9$ points. Encoding B also finds all $345$ structured instances with $m+8$ points
unsatisfiable. SAT enumeration uses PySAT~\cite{pysat}.

\begin{table}[ht]
\centering
\small\begin{tabular}{@{}llr@{}}
\toprule
bound & instances & proofs\\
\midrule
$f(5,6)<12$, $f(5,7)<13$ & encoding A, 1 instance each & $2$\\
$f(6,6)<13$, $f(6,7)<14$ & encoding A, 1 instance each & $2$\\
$f(6,8)<16$ & encoding A, split by the first column & $64$\\
$f(7,7)<15$ & encoding A, split by the first column & $128$\\
$f(7,8)<17$ & encoding A, split by the first column & $128$\\
$f(6,m)<m+9$ for $8\le m\le14$ & structured, 10 row graphs for each width & $70$\\
$f(6,m)<m+8$ for $9\le m\le11$ & structured, 115 row graphs for each width & $345$\\
\bottomrule
\end{tabular}
\caption{The unsatisfiable instances.}
\label{tab:instances}
\end{table}

The $72$ and $1440$ candidates with $m+4$ points on the $4\times4$ and $4\times5$ grids are
searched by \nolinkurl{verify_narrow.py}, which finds 2 rectangles with the same
presentation in each of them.

\section{Explicit blackouts}

The blackouts and test boards of Section 6 of the paper are checked from the definition by
the scripts in \nolinkurl{jis/scripts} and by
\nolinkurl{papers/extremal/scripts/check_constructions.py}. The second script lists the
tilted rectangles through their leftmost corner, compares all presentations with at most 2
points, and looks for an axis-aligned rectangle with the same presentation. It also checks
the explicit construction with spacing $W_h+1$ for $3\le h\le8$, and blackouts built as in
Section 5 for $4\le h\le14$ and $h=16$. In these, each double gets the leftmost column that
is not bad, where the conditions (B1), (B2), and (B3) are tested as defined. At $h=16$, all 3
blackouts built in this way are valid, and all 5 blackouts built with (B1) alone are not
valid. For $h=20$ and $h=24$ the script only computes the columns and prints the width. The
script runs in about 4 minutes.

\section{Formal proofs}

The Lean files are in \nolinkurl{lean/KageyProblems/Problem001}. They contain formal proofs
of the lemma on presentations, of the lemma that 3 corners determine a rectangle, of the
span bound, of the row bound, of the finite collar theorem, and of the alternating seas
proposition under the assumption that the doubles are $2W_h+2$ columns apart
(\nolinkurl{GeneralLocal.lean} and \nolinkurl{GeneralTheorem.lean}). This gives
$f(h,m)=m+\lfloor h^2/4\rfloor$ for every $h\ge3$ and
$m\ge q_h+(q_h-1)(2W_h+1)$. Validity tests of $S_0$, $S_8$, the $6\times8$ and $6\times15$
blackouts, and the $5\times13$ and $6\times23$ test boards are checked by the Lean kernel.
Together with the formal row bound and collar theorem, this proves $f(5,m)=m+6$ for
$m\ge8$, $f(6,m)=m+9$ for $m\ge15$, $f(5,5)\ge10$, and $f(6,8)\ge15$ in Lean.

Not formalized are the spacing $W_h+1$, the threshold theorem and its corollary, the
equality structure theorem, the lemma on 6 rows and $m+8$ points, the blackouts on 7 rows,
and Lemma~\ref{lem:enc}.

\begin{thebibliography}{9}
\bibitem{bailleux} O. Bailleux and Y. Boufkhad. Efficient CNF encoding of Boolean
cardinality constraints. In \emph{Principles and Practice of Constraint Programming -- CP
2003}, Lecture Notes in Comput. Sci. 2833, pages 108--122. Springer, 2003.
\bibitem{cadical} A. Biere, T. Faller, K. Fazekas, M. Fleury, N. Froleyks, and F. Pollitt.
CaDiCaL 2.0. In \emph{Computer Aided Verification, CAV 2024}, Lecture Notes in Comput. Sci.
14681, pages 133--152. Springer, 2024.
\bibitem{cruzfilipe} L. Cruz-Filipe, M. J. H. Heule, W. A. Hunt, Jr., M. Kaufmann, and P.
Schneider-Kamp. Efficient certified RAT verification. In \emph{Automated Deduction -- CADE
26}, Lecture Notes in Comput. Sci. 10395, pages 220--236. Springer, 2017.
\bibitem{pysat} A. Ignatiev, A. Morgado, and J. Marques-Silva. PySAT: a Python toolkit for
prototyping with SAT oracles. In \emph{Theory and Applications of Satisfiability Testing --
SAT 2018}, Lecture Notes in Comput. Sci. 10929, pages 428--437. Springer, 2018.
\bibitem{sinz} C. Sinz. Towards an optimal CNF encoding of Boolean cardinality
constraints. In \emph{Principles and Practice of Constraint Programming -- CP 2005},
Lecture Notes in Comput. Sci. 3709, pages 827--831. Springer, 2005.
\bibitem{drattrim} N. Wetzler, M. J. H. Heule, and W. A. Hunt, Jr. DRAT-trim: efficient
checking and trimming using expressive clausal proofs. In \emph{Theory and Applications of
Satisfiability Testing -- SAT 2014}, Lecture Notes in Comput. Sci. 8561, pages 422--429.
Springer, 2014.
\end{thebibliography}

\end{document}
